%% 5-0 序章　スピードメーターは何を計算している？（第11回：6ページ）
\makeatletter
\gdef\bk@partlabel{第5巻}\gdef\bk@partname{変化をつかむ}
\makeatother
\bkchapter{5}{0}{スピードメーターは\\何を計算している？}{\bklevel{2}}{スピードメーターは何を計算している？}

\noindent 車のスピードメーターの針が「50」を指しています。時速\qty{50}{km}です。でも、この車は1時間も走っていません。走りはじめて、まだ10秒です。

\begin{bkfigure}
\begin{tikzpicture}[x=1mm, y=1mm]
  % スピードメーター（0〜180 km/h）
  \fill[bkInk!88] (0,0) circle (34);
  \fill[bkInk!75] (0,0) circle (31);
  \foreach \v in {0,10,...,180} {
    \pgfmathsetmacro\ang{210-\v*4/3}
    \pgfmathtruncatemacro\isbig{mod(\v,20)==0}
    \ifnum\isbig=1
      \draw[white, line width=1.2pt] (\ang:24.5) -- (\ang:29.5);
      \node[white, font=\sffamily\bfseries\footnotesize] at (\ang:19.5) {\v};
    \else
      \draw[white, line width=0.6pt] (\ang:26.5) -- (\ang:29.5);
    \fi}
  \pgfmathsetmacro\need{210-50*4/3}
  \draw[bkOrange, line width=2.2pt, line cap=round] (0,0) -- (\need:27);
  \fill[bkOrange!70!black] (0,0) circle (3);
  \node[white, font=\sffamily\footnotesize] at (0,-12) {km/h};
  % 説明
  \node[anchor=west, text width=58mm, align=left, font=\small] at (42,11)
    {針は、いまの速さ「時速\qty{50}{km}」を指している。\\[2pt]
     でも、時速\qty{50}{km}は「1時間に\qty{50}{km}進む速さ」。この車は、まだ1時間走っていない。};
  \node[anchor=west, text width=58mm, align=left, font=\sffamily\bfseries\small, text=bkDeep] at (42,-18)
    {走っていない1時間の距離を、\\メーターはどうやって知ったのだろう？};
\end{tikzpicture}
\end{bkfigure}

\bksection{doguu}{「時速\qty{50}{km}」は、何を言っているの？}

時速\qty{50}{km}は「1時間あたり\qty{50}{km}」という意味で、\qty{50}{km/h}と書きます（h は hour の頭の文字で「時間」）。でも、この車は1時間走ったわけではありません。メーターが言っているのは、次のことです。

\begin{taisetsu}
\textbf{もし、いまの速さのまま1時間走りつづけたら、\qty{50}{km}進む。}
\end{taisetsu}

「いまの速さのまま」なら、1時間より短い時間でも、進む距離が計算できます。30分は1時間の半分なので、進む距離も半分の\qty{25}{km}です。1分では、1時間が60分（\qty{60}{min/h}、min は minute の頭の文字で「分」）なので、
\[ 50\,\mathrm{km/h}\div60\,\mathrm{min/h}=0.833\cdots\,\mathrm{km/min} \]
h どうしが約分されて消え、「1分あたり約\qty{0.83}{km}」（約\qty{833}{m}）になります。

\begin{bkfigure}
\begin{tikzpicture}[x=1mm, y=1mm]
  % 道路と、いまの車。点線の車は「いまの速さのまま」走ったときの位置（はばは目安で、長さの比ではない）
  \fill[bkGray!18] (-4,-2.5) rectangle (128,0);
  \draw[bkGray, line width=0.8pt] (-4,0) -- (128,0);
  \bkcar{(0,0)}{0.8}{bkOrange}
  \node[anchor=south, font=\sffamily\scriptsize, text=bkOrange!80!black] at (9.6,14.5) {いま};
  \draw[bk arrow, bkBlue, dash pattern=on 3pt off 2pt] (26,7) -- (124,7);
  \node[anchor=south, font=\sffamily\footnotesize, text=bkBlue] at (75,8.5) {いまの速さ（時速\qty{50}{km}）のまま、走りつづけたら……};
  \foreach \x/\lab in {40/1秒後, 64/10秒後, 88/1分後, 112/1時間後} {
    \draw[bkGray, line width=0.7pt] (\x,0) -- (\x,-4);
    \node[anchor=north, font=\sffamily\scriptsize] at (\x,-4.8) {\lab};}
\end{tikzpicture}
\end{bkfigure}

\begin{bkcenter}
\begin{bktable}{colspec={Q[l,m,wd=28mm] *{5}{Q[c,m,wd=17mm]}}, rowsep=1.5pt, colsep=3pt}
走りつづけた時間 & 1秒 & 10秒 & 1分 & 30分 & 1時間 \\
進む距離 & 約\qty{14}{m} & 約\qty{139}{m} & 約\qty{833}{m} & \qty{25}{km} & \qty{50}{km} \\
\end{bktable}
\bktabcaption{時速\qty{50}{km}のまま走ると、どれだけ進むか}
\end{bkcenter}

1秒に進む距離は、単位つきの式で計算します（2-10）。\qty{1}{km}は\qty{1000}{m}なので「1 km あたり 1000 m」で \qty{1000}{m/km}、1時間は $60\times60=3600$ 秒なので「1時間あたり3600秒」で \qty{3600}{s/h} です（s は second の頭の文字で「秒」）。
\[ 50\,\mathrm{km/h}\times1000\,\mathrm{m/km}\div3600\,\mathrm{s/h}=13.88\cdots\,\mathrm{m/s} \]
km どうし、h どうしが約分されて消え、m/s が残ります。だから1秒では約\qty{14}{m}進みます。くわしい計算のしかたは、5-1でもう一度ていねいに見ます。

%% ---------------- p.2 瞬間の速さ ----------------
\bksection{doguu}{「いまの速さ」は、どうやって計算する？}

ここで、こまったことが起きます。2-3で学んだ速さの式は、
\[ \text{速さ}=\text{進んだ道のり}\div\text{かかった時間} \]
でした。この式は、「10秒のあいだに\qty{100}{m}進んだ」のように、\textbf{ある時間のあいだ}の速さです。このあいだに速くなったり、おそくなったりしても、ならした値になります。だから\textbf{平均の速さ}といいます。\bkidx{はやさ}{速さ}

スピードメーターが表示しているのは、「いま、この瞬間」の速さです。瞬間とは、時間の長さが0の「時刻の1点」のことです。時間の長さが0だと、進んだ道のりも0です。すると、式は
\[ 0\div0 \]
になってしまいます。0でわる計算はできません（1-14）。

\begin{bkfigure}
\begin{tikzpicture}[x=1mm, y=1mm]
  % 時間の線。上：あいだ（平均の速さ）、下：瞬間（時刻の1点）
  \draw[bk arrow, bkInk, line width=0.8pt] (0,14) -- (122,14);
  \node[anchor=west, font=\sffamily\scriptsize] at (122.5,14) {時間};
  \fill[bkBlue!25] (24,11.8) rectangle (84,16.2);
  \draw[bkBlue, line width=0.8pt] (24,11.8) rectangle (84,16.2);
  \node[anchor=south, font=\sffamily\footnotesize, text=bkBlue] at (54,17.5) {あいだ（たとえば10秒）};
  \node[anchor=north, font=\sffamily\footnotesize, text=bkInk, align=center] at (54,10) {進んだ道のり÷かかった時間 で計算できる（平均の速さ）};
  \draw[bk arrow, bkInk, line width=0.8pt] (0,-6) -- (122,-6);
  \node[anchor=west, font=\sffamily\scriptsize] at (122.5,-6) {時間};
  \fill[bkOrange] (54,-6) circle (1.6);
  \node[anchor=south, font=\sffamily\footnotesize, text=bkOrange!80!black] at (54,-3) {瞬間（時刻の1点。長さ0）};
  \node[anchor=north, font=\sffamily\footnotesize, text=bkInk] at (54,-9.5) {道のり0÷時間0 になって、このままでは計算できない};
\end{tikzpicture}
\end{bkfigure}

それでもメーターは、「いまの速さ」を表示しています。どうやっているのでしょう。答えは5-1にあります。ヒントは、「0にする」のではなく、「とても短いあいだ」を考えることです。

この「瞬間の変化の速さ」を計算する道具が\textbf{微分}（びぶん）です。そして反対に、速さから進んだ距離を求める道具が\textbf{積分}（せきぶん）です。この巻では、この2つの道具を、表計算とグラフを使って、少しずつ組み立てていきます。

\begin{taisetsu}
\textbf{微分}：いま、どれだけの速さで変わっているか（瞬間の変化の速さ）\\
\textbf{積分}：少しずつの変化をたし合わせて、全部でどれだけになったか
\end{taisetsu}

\bkfukidashi{mina}{0÷0 はできないのに、メーターは計算している。ふしぎだね。}

%% ---------------- p.3 おふろで見る微分と積分 ----------------
\bksection{doguu}{おふろの水で見る、微分と積分}

微分と積分は、車の速さだけの道具ではありません。もう1つ、おふろに水をためる場面で考えてみましょう。

じゃ口を少しずつ開けながら、おふろに水を入れました。1分ごとに、たまった水の量を調べると、次のようになりました。

\begin{bkfigure}
\begin{tikzpicture}[x=1mm, y=1mm]
  % 左：浴そう。1分ごとに入った水を、色を変えて積み重ねる（6, 9, 12, 12, 12 L。高さは 1 L＝0.6 mm）
  \def\yy{0}
  \foreach \v/\c/\k in {6/bkBlue!18/1, 9/bkBlue!30/2, 12/bkBlue!42/3, 12/bkBlue!54/4, 12/bkBlue!66/5} {
    \pgfmathsetmacro\ya{\yy}\pgfmathsetmacro\yb{\yy+\v*0.6}
    \fill[\c] (2,\ya) rectangle (34,\yb);
    \node[font=\sffamily\scriptsize, text=bkInk] at (18,{(\ya+\yb)/2}) {\k 分めに \v\,L};
    \xdef\yy{\yb}}
  \draw[bkInk, line width=1.1pt, rounded corners=1.5mm] (0,38) -- (0,-1.2) -- (36,-1.2) -- (36,38);
  % じゃ口（右から浴そうの上へのびる管と、下向きの口）
  \fill[bkGray] (24,44) rectangle (40,47);
  \fill[bkGray] (24,40) rectangle (27.5,47);
  \draw[bkBlue!60, line width=1.2pt] (25.75,39.5) -- (25.75,31.5);
  \node[anchor=north, font=\sffamily\footnotesize] at (18,-3) {5分で\qty{51}{L}たまった};
  % 右：表。ふえた量は、となりの列のあいだに書く。2〜3分の1組だけ、ひき算の線をかく
  \begin{scope}[shift={(46,0)}]
    \foreach \t [count=\i from 0] in {0,1,2,3,4,5} {\node[font=\sffamily\small] at (26+\i*11,36) {\t};}
    \foreach \v [count=\i from 0] in {0,6,15,27,39,51} {\node[font=\sffamily\small] (w\i) at (26+\i*11,26) {\v};}
    \foreach \v [count=\i from 0] in {6,9,12,12,12} {\node[font=\sffamily\small, text=bkOrange!80!black] (d\i) at (31.5+\i*11,6) {\v};}
    \node[anchor=east, font=\sffamily\scriptsize, align=right] at (20,36) {時刻（分）};
    \node[anchor=east, font=\sffamily\scriptsize, align=right] at (20,26) {たまった量（L）};
    \node[anchor=east, font=\sffamily\scriptsize, align=right] at (20,6) {1分ごとに\\ふえた量（L）};
    \draw[bkLine, line width=0.5pt] (21,31) -- (86,31) (21,21) -- (86,21) (21,11) -- (86,11);
    \draw[bkOrange!80!black, line width=0.6pt] (w2.south) -- (d2.north) (w3.south) -- (d2.north);
    \node[anchor=west, font=\sffamily\scriptsize, text=bkOrange!80!black, inner sep=0.8pt] at (59.5,15.5) {$27-15=12$};
  \end{scope}
\end{tikzpicture}
\bkfigcaption{左：1分ごとに入った水が、積み重なってたまる。右：たまった量と、1分ごとにふえた量。}
\end{bkfigure}

\textbf{ひき算で、ふえ方がわかる}\par
たまった量の表から、となりどうしのひき算（あと${}-{}$まえ、1-5・2-3）をすると、1分ごとにふえた量がわかります。たとえば、2分から3分までは $27-15=12$ で、\qty{12}{L}ふえました。これは「1分あたり\qty{12}{L}」という、\textbf{水がたまる速さ}です。この「ふえ方（変化の速さ）を求める」考えを、とても細かい時間まで進めたものが\textbf{微分}です。

\textbf{たし算で、全部の量がわかる}\par
反対に、1分ごとにふえた量をはじめから順にたしていくと、たまった量がわかります。
\[ 6+9+12+12+12=51 \]
で、5分で\qty{51}{L}です。この「少しずつの量をたし合わせて全部を求める」考えを、とても細かい時間まで進めたものが\textbf{積分}です。

ここでは1分ごとに区切りました。でも、ほんとうは1分のあいだにも、水の入り方は少しずつ変わっています。区切りを0.1分、0.01分……と、どんどん細かくしていくと、微分と積分になります。微分はひき算のなかま、積分はたし算のなかまで、この2つはおたがいに逆の計算です（5-11）。

%% ---------------- p.4 この巻が通る路線 ----------------
\bksection{doguu}{この巻が通る路線と駅}

この巻の章は、4本の路線の駅です。前の巻で乗った路線が、この巻で「微分」と「積分」という大きな駅に着きます。

\begin{bkfigure}
\begin{tikzpicture}[x=1mm, y=1mm]
  \foreach \r/\y/\nm/\stA/\stB/\stC in {
      1/60/変化/{1-5・2-3・3-5}/{5-1〜5-8・5-17・5-18}/{6-1（微分方程式）},
      2/42/たし合わせ/{2-3・4-5}/{5-9〜5-13・5-16}/{6-3・8-11},
      3/24/倍々/{2-7・3-11}/{5-14}/{6-4・9-5},
      4/6/回転と波/{4-2・4-8}/{5-15}/{6-10・9-1}} {
    \draw[bkRoute\r, line width=3pt, line cap=round] (8,\y) -- (132,\y);
    \node[inner sep=0pt] at (4,\y) {\bkrouteicon[6.5mm]{\r}};
    \node[anchor=south west, font=\sffamily\bfseries\small, text=bkInk] at (10,\y+1.8) {\nm};
    \fill[white] (34,\y) circle (2.2);\draw[bkRoute\r, line width=1.1pt] (34,\y) circle (2.2);
    \node[anchor=north, font=\sffamily\scriptsize, text=bkGray] at (34,\y-3) {\stA};
    \fill[white, rounded corners=1.5mm] (56,\y-4.2) rectangle (104,\y+4.2);
    \draw[bkRoute\r, line width=1.6pt, rounded corners=1.5mm] (56,\y-4.2) rectangle (104,\y+4.2);
    \node[font=\sffamily\bfseries\scriptsize, text=bkInk] at (80,\y) {\stB};
    \fill[white] (124,\y) circle (2.2);\draw[bkRoute\r, line width=1.1pt] (124,\y) circle (2.2);
    \node[anchor=north, font=\sffamily\scriptsize, text=bkGray] at (124,\y-3) {\stC};}
  \node[font=\sffamily\bfseries\small, text=bkGray] at (34,72) {前の巻};
  \node[font=\sffamily\bfseries\small, text=bkDeep] at (80,72) {この巻（第5巻）};
  \node[font=\sffamily\bfseries\small, text=bkGray] at (124,72) {この先};
\end{tikzpicture}
\bkfigcaption{第5巻を通る4本の路線。四角の中がこの巻の駅。}
\end{bkfigure}

\begin{bkcenter}
\begin{bktable}{colspec={Q[c,m,wd=8mm] Q[l,m,wd=30mm] Q[l,m,wd=82mm]}, rowsep=2pt}
 & 路線 & この巻で着く駅 \\
\bkrouteicon{1} & 変化 & ひき算（1-5）→ 傾き（3-5）→ \textbf{微分}。瞬間の変化の速さ \\
\bkrouteicon{2} & たし合わせ & かけ算の面積（2-3）→ \textbf{積分}。細い長方形をたし合わせる \\
\bkrouteicon{3} & 倍々 & 指数関数（3-11）→ \textbf{e}。微分しても形が変わらない関数 \\
\bkrouteicon{4} & 回転と波 & 三角関数（4-8）→ \textbf{sin・cosの微分}。波は微分すると1/4周期ずれる \\
\end{bktable}
\end{bkcenter}

変化の路線（\bkrouteicon[0.9em]{1}）とたし合わせの路線（\bkrouteicon[0.9em]{2}）は、5-11で1つにつながります。おふろの水で見た「ひき算とたし算は逆の計算」が、「積分は微分の逆」という、この巻でいちばん大事な発見になります。

%% ---------------- p.5 この巻でできるようになること ----------------
\bksection{goal}{この巻でできるようになること}

\begin{bkfigure}
\begin{tikzpicture}[x=1mm, y=1mm,
  q/.style={draw=bkVol, line width=1.2pt, fill=white, rounded corners=2.5mm, minimum width=34mm, minimum height=17mm, align=center, font=\sffamily\small}]
  \node[q] (x) at (0,0) {\textbf{位置} $x$\\[1pt]{\footnotesize どこにいる？}};
  \node[q] (v) at (52,0) {\textbf{速度} $v$\\[1pt]{\footnotesize どれだけ速く動く？}};
  \node[q] (a) at (104,0) {\textbf{加速度} $a$\\[1pt]{\footnotesize 速さがどう変わる？}};
  \draw[bk arrow, bkOrange] ([yshift=3mm]x.east) -- node[above, font=\sffamily\bfseries\small] {微分} ([yshift=3mm]v.west);
  \draw[bk arrow, bkOrange] ([yshift=3mm]v.east) -- node[above, font=\sffamily\bfseries\small] {微分} ([yshift=3mm]a.west);
  \draw[bk arrow, bkBlue] ([yshift=-3mm]v.west) -- node[below, font=\sffamily\bfseries\small] {積分} ([yshift=-3mm]x.east);
  \draw[bk arrow, bkBlue] ([yshift=-3mm]a.west) -- node[below, font=\sffamily\bfseries\small] {積分} ([yshift=-3mm]v.east);
  \node[font=\footnotesize, text=bkGray] at (26,-15) {5-9〜5-13};
  \node[font=\footnotesize, text=bkGray] at (26,15) {5-1〜5-5};
  \node[font=\footnotesize, text=bkGray] at (78,15) {5-5};
  \node[font=\footnotesize, text=bkGray] at (78,-15) {5-13};
\end{tikzpicture}
\bkfigcaption{微分は右へ、積分は左へ。2つは逆向きの道具。}
\end{bkfigure}

位置・速度・加速度の3つは、「微分」で右へ、「積分」で左へ移れます。たとえば、位置が時間とともにどう変わるかがわかれば、微分して速度がわかります。反対に、速度がわかれば、積分して位置がわかります。おふろの水なら、「たまった量」が位置、「たまる速さ」が速度にあたります。

\begin{bkcenter}
\begin{bktable}{colspec={Q[l,m,wd=20mm] Q[l,m,wd=58mm] Q[l,m,wd=42mm]}, rowsep=2pt}
部 & できるようになること & 物理・工学の例 \\
第1部 微分 & 瞬間の速さを計算する。グラフの傾きを式で求める。いちばん高い所を見つける & スピードメーター、投げ上げたボールの最高点 \\
第2部 積分 & 速さのグラフの面積から進んだ距離を求める。たまった量を計算する & ばねの仕事、水の量、電気の使用量 \\
第3部 特別な関数 & e・sin・cos を微分・積分する。むずかしい関数を多項式で近似する & 放射線の減り方、ばねの振動、振り子 \\
\end{bktable}
\end{bkcenter}

\begin{column}{なぜ、微分と積分を学ぶの？}
物理の法則の多くは、「いまの状態が、次の瞬間の変化を決める」という形で書かれています。たとえば「落ちる物の速さは、1秒ごとに\qty{9.8}{m/s}ずつ増える」。この「変化の速さ」を書くことばが微分で、変化を積み重ねて全体を計算する道具が積分です。この巻の道具がそろうと、次の第6巻で「未来を計算する」微分方程式に進めます。\bkidx{びぶんほうていしき}{微分方程式}
\end{column}

%% ---------------- p.6 この巻の読み方 ----------------
\bksection{reidai}{この巻の読み方}

\bkfukidashi{sora}{微分って、むずかしそう。でも「平均の速さ」なら2-3でやったよ。}
\bkfukidashi{delta}{この巻では、まず表計算で数を出して、きまりを見つけてから式にするよ。手を動かせば、ちゃんとわかるよ。}

\bksubsection{1}{章のはじめの「思い出そう」}

この巻の章には、「思い出そう」の囲みがあります。その章で使う、前の巻で学んだことをまとめたものです。札の番号の章にもどれば、くわしく書いてあります。
\begin{omoidasou}
\bkrecall{2-3}{速さ＝進んだ道のり÷かかった時間。「変わった量」は \mbox{$\Delta$（デルタ）}で書く。}
\bkrecall{3-5}{グラフの傾き＝$\dfrac{y\text{の変化量}}{x\text{の変化量}}=\dfrac{\Delta y}{\Delta x}$}
\end{omoidasou}

\bksubsection{2}{表計算とノートを用意しよう}

表計算（2-18）をたくさん使います。「やってみよう」の表は、コンピュータで作ってみよう。問題はノートに解きます。電卓も使ってかまいません。

\bksubsection{3}{式の変形は、1行ずつ}

式の変形は1行ずつ書いてあります。どの行も、前の行に「展開する」「両辺を同じ数でわる」などのきまり（2-14・3-8）を1つ使っただけです。わからない行は、前の行とくらべてみよう。

\bksubsection{4}{単位の換算は、単位つきの式で}

時速と秒速、分と秒のように単位をなおすときは、\qty{1000}{m/km}（1 km あたり 1000 m）、\qty{60}{s/min}（1分あたり60秒）、\qty{3600}{s/h}（1時間あたり3600秒）のような「単位つきの数」をかけたりわったりします（2-10）。同じ単位どうしは約分されて消え、ほしい単位だけが残ります。単位が正しく残ったかどうかで、計算のまちがいにも気づけます。

\bksubsection{5}{わからなくなったら}

はじめての考えは、1回読んですぐにわからなくても大丈夫です。例題を、ノートに1行ずつ写しながら、自分の手で計算してみよう。どの行でわからなくなったかが見つかれば、その行の「（2-14）」のような番号の章にもどって確かめればよいのです。
